% !TeX root = ../泛函分析笔记.tex
% 由主文档通过 \input 载入；请编译主文档。
\setcounter{chapter}{1} % 保留原书章号。
\chapter{拓扑与集合论工具}
\label{app:fa-topology}

度量空间足以处理主线中的大多数例子，但一般弱紧性涉及任意坐标乘积，必须补充少量不依赖度量的语言。本附录区分定义、可直接验证的事实与作为先修知识引用的定理。进一步论证可参阅\cite{teschl-fa}的集合论及拓扑附录。

\section{拓扑、紧性与乘积}
集合 $X$ 上的拓扑 $\tau$ 是一族子集，包含 $\varnothing,X$，对任意并和有限交封闭；其成员称开集。映射连续是指开集的逆像开。任意两个不同点都有不交邻域时，空间称为 Hausdorff 空间。

紧性的定义仍是每个开覆盖有有限子覆盖。紧空间的闭子集紧，连续像紧；Hausdorff 空间中的紧子集闭。最后一个事实可如下证明：对紧集外的点 $x$，分别把它与紧集中每一点用不交邻域分开，再对紧集选有限子覆盖；相应的有限个 $x$ 邻域之交就是与紧集不交的邻域。

乘积 $\prod_{i\in I}X_i$ 的基本开集只限制有限个坐标：在有限指标集 $F\subset I$ 上要求 $x_i\in U_i$，其余坐标任意。本书引用\emph{Tychonoff 定理}：任意一族紧空间的乘积在乘积拓扑下紧。它是 Banach--Alaoglu 定理中唯一使用的任意乘积紧性结果。

\section{网与子网}
有向集是带有反身、传递关系的非空集合 $D$，且任意两个元素都有共同后继。\emph{网（net）}是以有向集为指标的族 $(x_\alpha)_{\alpha\in D}$；它趋于 $x$，是指每个 $x$ 的邻域最终包含全部网元素。

序列就是以自然数为指标的网。若 $x\in\overline A$，用 $x$ 的邻域按反向包含排序，在每个邻域 $U$ 中选择 $a_U\in U\cap A$，便得到趋于 $x$ 的网。反之，一个来自 $A$ 的网若趋于 $x$，则每个邻域与 $A$ 相交。因此网总能刻画闭包，而序列只有在额外条件下才能做到。

若 $\phi:E\to D$ 保序，且对每个 $\alpha\in D$，最终有 $\phi(\beta)\geq\alpha$，则 $(x_{\phi(\beta)})_{\beta\in E}$ 称为一个子网。一般紧性判别断言：拓扑空间紧当且仅当每个网都有收敛子网。这里引用该判别；它解释为什么一般弱星紧性不能直接改写成序列子列结论。

\section{Zorn 引理}
偏序关系除反身、传递外还满足反对称性。链是其中任意两个元素都可比较的子集；上界是大于等于该子集全部元素的元素；极大元是无法被严格扩大的元素，不要求它大于整个集合的所有元素。

本书采用 Zorn 引理：若非空偏序集的每条链都在该集合内有上界，则它存在极大元。这是与选择公理等价的集合论原则。在 Hahn--Banach 证明中，元素是定义在不同子空间上的受控延拓；链的上界由这些兼容映射之并给出。使用引理时必须验证并后的定义域是线性子空间、数值无冲突且控制不等式仍成立。

\section{半范数族与局部凸空间}
若向量空间的加法与数乘连续，称为拓扑向量空间（topological vector space）；若零点有凸邻域基，称为局部凸空间（locally convex space）。一族半范数 $(p_i)_{i\in I}$ 可用基本邻域
\[
 \{x:p_{i_1}(x)<\varepsilon,\ldots,p_{i_m}(x)<\varepsilon\}
\]
生成局部凸拓扑；若半范数族能分离非零向量，则拓扑为 Hausdorff。弱拓扑就是取 $p_f(x)=|f(x)|$ 得到的情形。

例如，对开集 $\Omega\subset\mathbb R^d$ 中的光滑函数，用 $p_{K,m}(u)=\max_{|\alpha|\leq m}\sup_{x\in K}|\partial^\alpha u(x)|$ 描述各阶导数在紧集 $K\subset\Omega$ 上一致收敛。这里 $\alpha=(\alpha_1,\ldots,\alpha_d)$ 为非负整数多重指标，$|\alpha|=\sum_j\alpha_j$。这种结构用于分布理论，不必强行由单个范数描述。

\section{习题}
\begin{enumerate}
 \item 直接证明坐标投影在乘积拓扑下连续。
 \item 用网的定义证明闭集包含来自该集合的所有收敛网的极限。
 \item 解释 Zorn 引理中的极大元与最大元的区别。
 \item 由半范数邻域验证弱拓扑下的加法和数乘连续。
\end{enumerate}

